\documentclass[11pt]{amsart}

\usepackage[T1]{fontenc}
\usepackage{lmodern}
\usepackage{microtype}
\usepackage{mathtools,amssymb,amsthm}
\usepackage[hidelinks]{hyperref}

\hypersetup{
  pdftitle={An order-49 partition refutes the heavy-chopped-rectangle interval
            conjecture}
}

\newtheorem{theorem}{Theorem}
\newtheorem{lemma}[theorem]{Lemma}
\newtheorem{proposition}[theorem]{Proposition}
\theoremstyle{definition}
\newtheorem{conjecture}[theorem]{Conjecture}
\theoremstyle{remark}
\newtheorem*{remark}{Remark}

\newcommand{\br}[1]{[#1]}
\newcommand{\sgv}{\operatorname{sg}}
\newcommand{\mex}{\operatorname{mex}}
\newcommand{\str}{\mathrm{st}}

\title[An order-49 counterexample to a heavy-interval conjecture]
{An Order-49 Partition Refutes the Heavy-Chopped-Rectangle
 Interval Conjecture of Gottlieb, Krnc and Mur\v{s}i\v{c}}

\date{}

\subjclass[2020]{91A46, 05A17, 05E10}
\keywords{combinatorial game, Sprague--Grundy value, integer partition,
Young's lattice, heavy position, counterexample}

\begin{document}

\begin{abstract}
In $1$-PNim a move deletes a positive number of rows, or a positive number of
columns, from the Young diagram of a single partition; a position is
\emph{heavy} when its Sprague--Grundy value equals the length of the longest
play from it. Gottlieb, Krnc and Mur\v{s}i\v{c} conjectured that whenever the
rectangle $\br{(a+1)^{b+1}}$ is heavy, every $\lambda$ with
$\br{a+1,a,\dots,a-b+1}\le\lambda\le\br{(a+1)^{b+1}}$ in Young's lattice is
heavy. We refute this. For $(a,b)=(10,4)$ the rectangle $\br{11^5}$ is heavy,
yet $\lambda=\br{11,11,11,8,8}$, which lies strictly between the two
endpoints, has $\sgv(\lambda)=10$ where the conjecture demands
$a+b+1=15$. The certificate is a mex table over the $45$ options of $\lambda$,
printed here in full. No counterexample has order less than $49$, and exactly
one other has order $49$.
\end{abstract}

\maketitle

\section{The game and the conjecture}

Let $\lambda=\br{\lambda_1,\dots,\lambda_r}$ be a partition, written as a
weakly decreasing tuple of positive parts, and let
$\lambda'=\br{\lambda'_1,\dots,\lambda'_{\lambda_1}}$ be its conjugate. In the
game $1$-PNim of Gottlieb, Krnc and Mur\v{s}i\v{c}~\cite{GKM}, the moves from
$\lambda$ are
\begin{enumerate}
\item[(1)] $\lambda\to\br{\lambda_{i_1},\dots,\lambda_{i_\ell}}$ for a proper
  subsequence $i_1,\dots,i_\ell$ of $1,\dots,r$; and
\item[(2)] $\lambda\to\br{\lambda'_{i_1},\dots,\lambda'_{i_\ell}}'$ for a
  proper subsequence $i_1,\dots,i_\ell$ of $1,\dots,\lambda_1$,
\end{enumerate}
that is, a move deletes a positive number of rows, or a positive number of
columns, and the surviving rows (columns) merge. The empty partition is the
unique terminal position. We write $\sgv$ for the Sprague--Grundy value under
normal play, so that $\sgv(\lambda)=\mex\{\sgv(\mu):\lambda\to\mu\}$ and
$\sgv(\br{\,})=0$.

The longest play from a nonempty $\lambda$ has length
$L(\lambda):=\lambda_1+r-1$, whence $\sgv(\lambda)\le L(\lambda)$
\cite[Prop.~2.6]{GKM}; the position $\lambda$ is called \emph{heavy} when
equality holds. Two families of heavy positions are relevant here. First,
\begin{equation}\label{eq:rect}
  \sgv\bigl(\br{c^{\,r}}\bigr)=\bigl((r-1)\oplus(c-1)\bigr)+1
\end{equation}
for all positive $r,c$ \cite[Thm.~1.1]{GKM}, where $\oplus$ is the Nim-sum, so
$\br{c^{\,r}}$ is heavy exactly when $(r-1)\oplus(c-1)=(r-1)+(c-1)$, i.e.\
when $\binom{c+r-2}{r-1}$ is odd. Second, the staircase
$\str_{c,r}:=\br{c,c-1,\dots,c-r+1}$ is heavy for all $c\ge r\ge1$
\cite[Prop.~4.4]{GKM}. Young's lattice is the order in which
$\br{\lambda_1,\dots,\lambda_\ell}\le\br{\mu_1,\dots,\mu_m}$ means
$\ell\le m$ and $\lambda_j\le\mu_j$ for $j=1,\dots,\ell$.

The conjecture we refute observes that the two endpoints of a certain interval
of Young's lattice are heavy for the two reasons just quoted, and asks whether
everything between them is heavy as well. It carries no printed number: it
appears in the open-problems section of~\cite{GKM} under the source's own
label \texttt{conj:heavychoppedrect}, which is the only stable handle the
source provides, and we therefore quote it verbatim from the source file of the
e-print of \texttt{2506.04991v1}:

\begin{quote}\ttfamily\small
\begin{tabbing}
\verb|\begin{conj}\label{conj:heavychoppedrect}|\\
\verb|    Let $a,b\in\ZZ$ be such that $\br{(a+1)^{b+1}}$ is heavy.|\\
\verb|      %a\ge b$ and $a\oplus b=a+b$. |\\
\verb|    Then any $\lambda$ satisfying|\\
\verb|    $\br{a+1,a,\dots, a-b+1}\le\lambda\le \br{(a+1)^{b+1}}$ is heavy.|\\
\verb|    %that is, $\sgpnim(\lambda)=a+b+1$.|\\
\verb|\end{conj}|
\end{tabbing}
\end{quote}

\noindent In display, therefore:

\begin{conjecture}[Gottlieb--Krnc--Mur\v{s}i\v{c}]\label{conj:target}
Let $a,b\in\mathbb Z$ be such that $\br{(a+1)^{b+1}}$ is heavy. Then any
$\lambda$ satisfying
$\br{a+1,a,\dots,a-b+1}\le\lambda\le\br{(a+1)^{b+1}}$ is heavy.
\end{conjecture}

Three remarks fix the reading, and none of them creates an escape.
\emph{First}, two clauses of the hypothesis are commented out in the source and
so do not appear in the compiled paper: $a\ge b$ and $a\oplus b=a+b$. The
second is not a restriction at all: by~\eqref{eq:rect} the printed hypothesis
``$\br{(a+1)^{b+1}}$ is heavy'' says $\binom{a+b}{b}$ is odd, which by the
parity criterion for binomial coefficients \cite{Fine} is exactly
$a\wedge b=0$, i.e.\ $a\oplus b=a+b$. The first is a real restriction, and it
is needed for the lower endpoint to be a partition at all: if $b>a$ then
$\br{a+1,a,\dots,a-b+1}$ has nonpositive entries. The counterexample below has
$a=10\ge4=b$, so it refutes Conjecture~\ref{conj:target} under either reading.
\emph{Second}, both endpoints have first part $a+1$ and exactly $b+1$ rows, so
by Young's lattice every $\lambda$ in the interval has $\lambda_1=a+1$ and
$r=b+1$ as well; hence $L(\lambda)=a+b+1$ for every member, and ``$\lambda$ is
not heavy'' and ``$\sgv(\lambda)<a+b+1$'' are the same statement, which is what
the third commented line records.
\emph{Third}, the two endpoints really are heavy under the hypothesis: the
upper one by assumption, the lower one because it is $\str_{a+1,b+1}$.

\section{The counterexample}

\begin{theorem}\label{thm:main}
Let $(a,b)=(10,4)$ and $\lambda=\br{11,11,11,8,8}$, a partition of $49$. Then
\begin{enumerate}
\item $\br{11^5}$ is heavy, so the hypothesis of
  Conjecture~\ref{conj:target} holds at $(a,b)=(10,4)$;
\item $\br{11,10,9,8,7}\le\lambda\le\br{11^5}$ in Young's lattice, and
  $\lambda$ is neither endpoint; and
\item $\sgv(\lambda)=10$, whereas $L(\lambda)=a+b+1=15$.
\end{enumerate}
Consequently $\lambda$ is not heavy and Conjecture~\ref{conj:target} is false.
\end{theorem}

\begin{proof}
(1) $\binom{a+b}{b}=\binom{14}{4}=1001$ is odd, equivalently
$10\wedge4=0$; and~\eqref{eq:rect} gives
$\sgv(\br{11^5})=\bigl((5-1)\oplus(11-1)\bigr)+1=(4\oplus10)+1=14+1=15
=11+5-1=L(\br{11^5})$.

(2) The lower endpoint is $\str_{11,5}=\br{11,10,9,8,7}$, and
$11\le11$, $10\le11$, $9\le11$, $8\le8$, $7\le8$ on five rows; the upper
endpoint dominates $\lambda$ because every part of $\lambda$ is at most $11$
and both have five rows. Neither endpoint equals $\lambda$.

(3) We first list the options. The rows of $\lambda$ are $11,11,11,8,8$, so a
move of type~(1) keeps $i$ of the three $11$s and $j$ of the two $8$s with
$0\le i\le3$, $0\le j\le2$ and $(i,j)\ne(3,2)$: that is $12-1=11$ options
$\br{11^i,8^j}$. The conjugate is $\lambda'=\br{5^8,3^3}$, so a move of
type~(2) keeps $p$ of the eight $5$s and $q$ of the three $3$s with
$0\le p\le8$, $0\le q\le3$ and $(p,q)\ne(8,3)$, giving
$\br{5^p,3^q}'=\br{(p+q)^3,p^2}$ after deleting zero parts: that is
$36-1=35$ options, pairwise distinct because $(p,q)$ is recovered from
$(p+q,p)$. A partition occurring in both lists with a positive part would need
either five parts in the pattern $3+2$, which among the row options forces the
excluded $(i,j)=(3,2)$, or three equal parts $\br{q^3}$ with $q\le3$ matching
$\br{11^3}$ or $\br{8^3}$, which is impossible. So the two lists meet only in
$\br{\,}$ and
\[
  \#\{\mu:\lambda\to\mu\}=11+35-1=45 .
\]

Next, no option can carry the value $15$. Every option loses at least one row
or at least one column and gains neither, so $\mu_1+\ell(\mu)\le
\lambda_1+r-1=15$ for every option $\mu$, i.e.\ $L(\mu)\le14$; by
$\sgv\le L$ this gives $\sgv(\mu)\le14$. (This step is the source's own, inside
the proof of \cite[Prop.~2.6]{GKM}.) Hence
\begin{align*}
  \lambda\ \text{is heavy}
  &\iff \sgv(\lambda)=15\\
  &\iff \text{the options of }\lambda\text{ realise every value }0,1,\dots,14,
\end{align*}
and exhibiting one value in $\{0,\dots,14\}$ realised by no option suffices.
Table~\ref{tab:mex} lists all $45$ options with their Sprague--Grundy values.
The value set is
\[
  \{0,1,2,3,4,5,6,7,8,9\}\cup\{11,12,13,14\},
\]
so $10$ is realised by no option while $0,\dots,9$ all are, and
$\sgv(\lambda)=\mex=10$. Since $10<15$, $\lambda$ is not heavy.
\end{proof}

\begin{table}[htb]
\small
\begin{tabular}{llr@{\qquad}llr@{\qquad}llr}
$\mu$ & & $\sgv$ & $\mu$ & & $\sgv$ & $\mu$ & & $\sgv$\\
\hline
$\br{10,10,10,8,8}$ &        & $14$ & $\br{7,7,7,5,5}$ &        & $11$ & $\br{4,4,4,3,3}$ &        & $8$\\
$\br{10,10,10,7,7}$ &        & $14$ & $\br{6,6,6,6,6}$ & $\ast$ & $2$  & $\br{4,4,4,2,2}$ &        & $8$\\
$\br{9,9,9,8,8}$    &        & $13$ & $\br{11,11,8}$   &        & $13$ & $\br{8,8}$       & $\ast$ & $7$\\
$\br{9,9,9,7,7}$    &        & $13$ & $\br{7,7,7,4,4}$ &        & $11$ & $\br{3,3,3,3,3}$ & $\ast$ & $7$\\
$\br{11,11,11,8}$   &        & $14$ & $\br{6,6,6,5,5}$ &        & $2$  & $\br{4,4,4,1,1}$ &        & $8$\\
$\br{8,8,8,8,8}$    & $\ast$ & $4$  & $\br{11,8,8}$    &        & $13$ & $\br{3,3,3,2,2}$ &        & $7$\\
$\br{9,9,9,6,6}$    &        & $13$ & $\br{6,6,6,4,4}$ &        & $2$  & $\br{3,3,3,1,1}$ &        & $7$\\
$\br{8,8,8,7,7}$    &        & $12$ & $\br{5,5,5,5,5}$ & $\ast$ & $1$  & $\br{11}$        & $\ast$ & $11$\\
$\br{11,11,8,8}$    &        & $14$ & $\br{6,6,6,3,3}$ &        & $2$  & $\br{2,2,2,2,2}$ & $\ast$ & $6$\\
$\br{8,8,8,6,6}$    &        & $12$ & $\br{5,5,5,4,4}$ &        & $1$  & $\br{3,3,3}$     & $\ast$ & $1$\\
$\br{7,7,7,7,7}$    & $\ast$ & $3$  & $\br{11,11}$     & $\ast$ & $12$ & $\br{2,2,2,1,1}$ &        & $6$\\
$\br{8,8,8,5,5}$    &        & $12$ & $\br{5,5,5,3,3}$ &        & $1$  & $\br{8}$         & $\ast$ & $8$\\
$\br{7,7,7,6,6}$    &        & $4$  & $\br{4,4,4,4,4}$ & $\ast$ & $8$  & $\br{2,2,2}$     & $\ast$ & $4$\\
$\br{11,11,11}$     & $\ast$ & $9$  & $\br{5,5,5,2,2}$ &        & $9$  & $\br{1,1,1,1,1}$ & $\ast$ & $5$\\
$\br{11,8}$         & $\dag$ & $12$ & $\br{1,1,1}$     & $\ast$ & $3$  & $\br{\,}$        & $\ast$ & $0$\\
\end{tabular}
\medskip
\caption{All $45$ options of $\lambda=\br{11,11,11,8,8}$, listed in the order
of decreasing size within each column. The value $10$ does not occur, and
$0,\dots,9$ all do, so $\sgv(\lambda)=10$. A star marks the $17$ entries whose
value is forced without computation --- the $16$ rectangles, by~\eqref{eq:rect},
together with $\br{\,}$; a dagger marks the one further entry forced by the
two-row formula of \cite[Prop.~4.9]{GKM}. The $17$ starred values already
comprise $\{0,\dots,9,11,12\}$, so $\sgv(\lambda)\ge10$ needs nothing beyond a
Nim-sum; the decisive half --- that no option carries $10$ --- rests on the
remaining $27$ values.}
\label{tab:mex}
\end{table}

\section{Least order}

Call a pair $(a,b)$ with $a\ge b\ge0$ and $a\wedge b=0$ \emph{admissible}, so
that Conjecture~\ref{conj:target} applies to it.

\begin{proposition}\label{prop:least}
No member of the interval of an admissible cell of order less than $49$ fails
to be heavy, and exactly two of order $49$ do: $\br{11,11,11,8,8}$ in the cell
$(10,4)$, with $\sgv=10$ against a required $15$, and
$\br{9,9,7,7,6,5,4,2}$ in the cell $(8,7)$, with $\sgv=3$ against a required
$16$.
\end{proposition}

The search range is finite for the following reason. The least-order member of
the cell $(a,b)$ is its lower endpoint $\str_{a+1,b+1}$, of order
$(b+1)(2a+2-b)/2$, so a cell can hold a member of order at most $49$ only if
that quantity is at most $49$; this leaves exactly $78$ admissible cells,
namely all $49$ with $b=0$, then $b=1$ with $a$ even and $a\le24$, $b=2$ with
$a\in\{4,5,8,9,12,13,16\}$, $b=3$ with $a\in\{4,8,12\}$, $b=4$ with
$a\in\{8,9,10\}$, $b=5$ with $a=8$, $b=6$ with $a\in\{8,9\}$, and $b=7$ with
$a=8$ --- every admissible cell with $b\ge8$ has least order at least $117$.
The $49$ cells with $b=0$ need no work, their interval being the single point
$\br{a+1}$, heavy by hypothesis. Proposition~\ref{prop:least} is then the
result of enumerating, in each of the remaining $29$ cells, every member of
order at most $49$ and evaluating $\sgv$; unlike Theorem~\ref{thm:main} it
rests on that computation and not on a hand-checkable certificate.

Nothing here refutes any of the heaviness results proved in \cite{GKM}, and no
cell is claimed free of counterexamples above order $49$.

\section{Reproduction}

Everything in Theorem~\ref{thm:main} except the $27$ unstarred entries of
Table~\ref{tab:mex} can be checked with a pen: $\binom{14}{4}=1001$, the
Nim-sum $4\oplus10=14$, five comparisons of parts for the membership, the
combinatorial count $11+35-1=45$, and the closed form~\eqref{eq:rect} on the
$16$ rectangular options. Those checks alone give $\sgv(\lambda)\ge10$, since
the starred values already realise $0,\dots,9$. To close the argument by hand
one must additionally confirm that none of the $27$ remaining options has value
$10$; each of those is a mex over its own option set, and the whole recursion
below $\lambda$ visits $194$ distinct partitions.

The program \texttt{verify.py} accompanying this note re-derives every
computational claim above from the move rule alone, in exact integer
arithmetic, reading the objects as printed here; run it as
\texttt{python3 verify.py} with Python 3.9 or later and the standard library
only. It prints the full $45$-row table of Table~\ref{tab:mex}, the sweep
behind Proposition~\ref{prop:least}, and one line per check, and exits nonzero
unless every check passes. The recorded run \texttt{verify.output.txt} reports
$87$ passing checks and closes with a statement of what it does not cover, in
particular that the closed forms quoted above and the source's two appendix
tables are transcribed from the e-print rather than checked against it. Neither
file is needed in order to redo the calculation of Theorem~\ref{thm:main},
since every object it consumes is printed above.

\begin{thebibliography}{9}

\bibitem{GKM}
E.~Gottlieb, M.~Krnc, and P.~Mur\v{s}i\v{c},
\emph{Nim on integer partitions and hyperrectangles},
preprint, 5 June 2025.
\href{https://arxiv.org/abs/2506.04991v1}{arXiv:2506.04991v1}.
Numbered statements are cited as they appear in the compiled \texttt{v1}
preprint; each is also stated in full above.

\bibitem{Fine}
N.~J. Fine,
\emph{Binomial coefficients modulo a prime},
Amer. Math. Monthly \textbf{54} (1947), no.~10, 589--592.
\href{http://www.jstor.org/stable/2304500}{jstor:2304500}.

\end{thebibliography}

\end{document}
